\section{Introduction}

Parameter fine-tuning is either done through additive, selective, re-parametrized or hybrid methods. The first involves introducing additional parameters to the model which are then trained, the second identifies a selection of parameters to retrain, the third decomposes the parameters into a low-rank representation which is then trained, and the fourth involves a combination of the previous. By utilising the fine-control of the parameterisation of deep network through the spline theory, deep network tuning can be done architecturally in proposed method named curvature tuning.

\section{Smoothing Out Region Assignment}

Recall that region assignment in a deep network operators according to an operation of the form 
\begin{equation}\label{eq:region_assignment}
    \left[\mathbf{t}^{(\ell)}\right]_i=\argmax_{j=1,\dots,r^{(\ell)}}\left(\left\langle\left[A^{(\ell)}\right]_{i,j\cdot},\mathbf{z}^{(\ell-1)}\right\rangle+\left[B\right]_{i,j}\right),
\end{equation}
which in Section \ref{sec:vector_quantization} we viewed as vector quantization. This lead to the $\beta$-VQ inference formulation, which showed that region assignment can be probabilistically relaxed to smooth out the non-linearity of deep networks
\begin{equation}\label{eq:vq_curvature_tuning}
    \left[\mathbf{t}_{\beta}^{(\ell)}\right]_i=\frac{\exp\left(\frac{\left[\beta^{(\ell)}\right]_i}{1-\left[\beta^{(\ell)}\right]_i}\left(\left\langle\left[A^{(\ell)}\right]_{i,j,\cdot},\mathbf{z}^{(\ell-1)}\right\rangle+\left[B^{(\ell)}\right]_{i,j}\right)\right)}{\sum_{t=1}^{r^{(\ell)}}\exp\left(\frac{\left[\beta^{(\ell)}\right]_i}{1-\left[\beta^{(\ell)}\right]_i}\left(\left\langle\left[A^{(\ell)}\right]_{i,t,\cdot},\mathbf{z}^{(\ell-1)}\right\rangle+\left[B^{(\ell)}\right]_{i,t}\right)\right)}.
\end{equation}
In particular, recall Proposition \ref{prop:relu_beta_vq} which shows how we can manoeuvrer from the linear regime, to the ReLU activation operator and then to the sigmoid gated linear activation operator simply by varying the parameter $\beta^{(\ell)}$ from zero to one. This has the effect of smoothing out the non-linearity of the deep network.

However, there is an alternative perspective we can take to induce a similar smoothing effective. More specifically, we can use a smooth the maximum operation in \eqref{eq:region_assignment} with
\begin{equation}\label{eq:max_curvature_tuning}
    \left[\tilde{\mathbf{t}}_{\beta}^{(\ell)}\right]_i=\left(1-\beta^{(\ell)}\right)\log\left(\sum_{j=1}^{r^{(\ell)}}\exp\left(\frac{1}{1-\beta^{(\ell)}}\left(\left\langle\left[A^{(\ell)}\right]_{i,j,\cdot},\mathbf{z}^{(\ell-1)}\right\rangle\right)+\left[B^{(\ell)}\right]_{i,j}\right)\right).
\end{equation}
In this case, as $\beta^{(\ell)}$ goes from zero to one we manoeuvrer from the linear regime to the ReLU activation operator.

\section{Curvature Tuning}

Given a trained deep network, we can parametrise its behaviour with the parameter $\beta^{(\ell)}$ using \eqref{eq:vq_curvature_tuning} or \eqref{eq:max_curvature_tuning}. In particular, note that regardless of which parameterisation we utilise, changing the value of $\beta^{(\ell)}$ will cause the decision boundary to shift. For tuning the deep network, this behaviour is undesirable. However, using an average of these parameterisation negates these shifts.

Curvature tuning utilises this and proposes to replace the ReLU activation operators of a trained deep network with $$\sigma_{\beta}(x)=\frac{\sigma_{\text{sig}}\left(\frac{\beta x}{1-\beta}\right)x}{2}+\frac{\log\left(1+\exp\left(\frac{x}{1-\beta}\right)\right)(1-\beta)}{2},$$and then tune the parameter $\beta$ using a forward passes on the deep network and a test set.

\subsection{References}

Chapter \ref{ch:finetuning} from \citet{huCurvatureTuningProvable2025}.